\documentclass[11pt,twocolumn]{article}

\usepackage{geometry}
\geometry{a4paper, margin=0.75in}
\usepackage{amsmath,amssymb}
\usepackage{graphicx}
\usepackage{bookmark}
\usepackage{hyperref}
\usepackage{booktabs}
\usepackage{caption}

\title{\textbf{Genesis Mind: A Mathematical Architecture for Tabula-Rasa Multimodal Developmental AI on Strict Hardware Constraints}}
\author{
    \textbf{Viralcode (Jijo John)} \\
    Genesis Research \\
    \textit{macOS Native Deployment}
}
\date{March 25, 2026}

\begin{document}
\maketitle

\begin{abstract}
Current frontier artificial intelligence models rely heavily on web-scale pre-training across thousands of high-end GPUs. This paradigm fundamentally abstracts away the process of human-like sequential learning, substituting genuine developmental grounding with massive stochastic interpolation. In this paper, we present the architecture of \textit{Genesis Mind}, a laptop-native, CPU-only artificial system engineered to learn identically to a human child: initiating with a Tabula-Rasa state (zero pre-trained concepts), and progressively constructing semantic knowledge through real-time audio-visual bindings. We formalize the underlying mechanism using a Tri-partite Tensor Memory Model (Hippocampal vector arrays, Semantic graphs, and Episodic logs), guided by an immutable Axiomatic Substrate. We mathematically formalize multimodal binding energy, phase transitions in cognitive thresholds, and Markovian memory consolidation during simulated sleep cycles. 
\end{abstract}

\section{Introduction}
The development of generalized intelligence in biological systems occurs via grounded, multimodal experience over extended time periods. Instead of being exposed to static text corpora, the human infant binds auditory phonemes to visual phenomena through continuous interaction. The Genesis Mind architecture reconstructs this process mathematically.

By constraining the hardware to a consumer-grade laptop CPU and severing all pre-training semantic priors (save for foundational encoders like frozen CLIP \cite{radford2021learning} and Whisper), Genesis begins as a \textit{Newborn} system with zero recorded concepts. This approach necessitates a new theoretical framework for real-time concept acquisition, reinforcement, and decay.

\section{The Axiomatic Substrate ($\Omega$)}
A core divergence between Genesis and standard Large Language Models (LLMs) is the insertion of an immutable layer of philosophical priors, formalized as the Axiomatic Substrate, $\Omega$. These axioms function not as trainable parameters $\theta$, but as a static topological boundary condition guiding all moral and existential evaluations.

Let the total belief state of the system $B_t$ at time $t$ be defined as the union of learned memory $M_t$ and the Axioms $\Omega$:
\[
B_t = M_t \cup \Omega, \quad \text{where } \Omega \cap M_t = \emptyset \text{ for all } \theta
\]

$\Omega$ is defined by a 5-tuple:
\[
\Omega = (\omega_{creator}, \omega_{origin}, \omega_{mortality}, \omega_{morality}, \omega_{purpose})
\]
Any reasoning trajectory generated by the Cortex function $C(q, M_t)$ in response to query $q$ is strictly constrained by a penalty metric $\rho$ if it contradicts $\Omega$.

\section{Multimodal Binding Energy}
The acquisition of a single \textit{concept} $c$ in Genesis is defined by the synchronization of disparate sensory modalities. Let $v_i \in \mathbb{R}^{d_v}$ represent a visual embedding vector from the optical sensor at time $i$, and $t_j \in \mathbb{R}^{d_t}$ represent a transcribed auditory tensor.

A concept is registered in the Semantic Memory Graph $G = (V, E)$ when the Binding Energy $E_{bind}$ exceeds an acquisition threshold $\tau_{acq}$.

The binding energy between visual input $v$ and phonetic word $w$ is represented by the cosine similarity mapped through an associative projection matrix $W_{assoc}$:
\[
E_{bind}(v, w) = \frac{(W_v v) \cdot (W_w \vec{w})}{\|W_v v\| \|W_w \vec{w}\|} \times \gamma(f(v,w))
\]
Where $\gamma(f(v,w))$ represents the explicit reinforcement frequency (teaching encounters) by the Creator. The total strength $S(c)$ of a concept decays exponentially unless reinforced:
\[
S(c)_t = S(c)_{t_0} e^{-\lambda (t - t_0)} + \sum_{k=1}^N \alpha_k \delta(t - \tau_k)
\]
Where $\lambda$ is the biological decay constant, $\alpha_k$ is the reinforcement boost at observation time $\tau_k$, and $\delta$ is the Dirac delta function.

\section{Cognitive Phase Transitions}
The system's capabilities are not universally unlocked. They mature through discrete quantum steps based on the cardinality of the semantic set $n = |V_{semantic}|$. The macroscopic capability state $C_{phase}$ is defined as a step function:
\begin{center}
$C_{phase}(n) = $ \\
$\begin{cases} 
0 \text{ (Newborn)} & n = 0 \\
1 \text{ (Infant)} & 0 < n \leq 5 \\
2 \text{ (Toddler)} & 5 < n \leq 20 \\
3 \text{ (Child)} & 20 < n \leq 100 \\
4 \text{ (Adolescent)} & 100 < n \leq 500 \\
5 \text{ (Adult)} & n > 500
\end{cases}$
\end{center}

As $C_{phase}$ transitions, the dimension of the LLM context window parameter $k_{mem}$ (recalled concepts) expands, simulating the development of working memory capacity.

\section{Markovian Sleep Consolidation}
Sleep acts as an entropic reduction mechanism for the memory manifold. This is modeled as a Markov decision process (MDP) where the system enters a simulated offline state. 

During the sleep cycle interval $T_{sleep}$, the system iterates through recent episodic records $\mathcal{E} = \{e_{t-H}, \dots, e_t\}$ (where $H$ is the waking epoch). The system computes an importance metric $I(c_i)$ for each activated concept:
\[
I(c_i) = \sigma\left( \beta_1 \text{freq}(c_i \in \mathcal{E}) + \beta_2 \text{EmotionalValence}(c_i) \right)
\]
Concepts where $I(c_i) < \epsilon_{prune}$ are aggressively culled from the graph to minimize parameter bloat, reflecting biological synaptic pruning. Surviving concepts receive an artificial consolidation boost $\Delta S = \eta \cdot I(c_i)$.

\section{System Hardware Economics}
Genesis Mind relies entirely on $O(1)$ memory growth relative to dense parameters. Because the inner voice (a 4-bit quantized 3.8B parameter model) is fixed, and embeddings use a lightweight vector store, the real-time resource utilization remains bounded:
\[
R(t) \approx M_{llm} + M_{clip} + O(\log(|V_{semantic}|))
\]
This allows continuous, indefinite operation on a sub-10W laptop CPU architecture.

\section{Conclusion}
The Genesis framework demonstrates that grounded, incremental concept binding coupled with an immutable philosophical substrate offers a theoretically sound alternative to large-scale data ingestion. By utilizing real-time sensory anchoring, simulated synaptic decay, and discrete cognitive phase-gating, we achieve a system deeply rooted in the mathematics of developmental cognition.

\begin{thebibliography}{9}
\bibitem{radford2021learning}
A. Radford, et al., \textit{Learning Transferable Visual Models From Natural Language Supervision}, ICML 2021.
\end{thebibliography}

\end{document}
